\documentclass[11pt]{article}

\usepackage[utf8]{inputenc}
\usepackage[T1]{fontenc}
\usepackage{lmodern}
\usepackage{geometry}
\usepackage{amsmath,amssymb,amsfonts,amsthm,mathtools}
\usepackage{bm}
\usepackage{enumitem}
\usepackage{microtype}
\usepackage{hyperref}
\usepackage{titlesec}
\usepackage{booktabs}
\usepackage{array}
\usepackage{setspace}
\usepackage{mathrsfs}
\usepackage{physics}

\geometry{margin=1in}
\hypersetup{colorlinks=true, linkcolor=black, urlcolor=blue, citecolor=black}
\setlist[enumerate]{topsep=0.4em,itemsep=0.35em}
\setlist[itemize]{topsep=0.3em,itemsep=0.25em}
\titleformat{\section}{\large\bfseries}{\thesection.}{0.5em}{}
\titleformat{\subsection}{\normalsize\bfseries}{\thesubsection.}{0.5em}{}
\setstretch{1.06}

\newcommand{\Aether}{\AE{}ther}
\newcommand{\Aflow}{\AE{}ther-flow}
\newcommand{\TheoryName}{\Aflow{} Interpretation of Relativity}
\newcommand{\TheoryProgram}{\Aether{} / \Aflow{} framework}

\newcommand{\CompletedMetric}{g^{\sharp}}
\newcommand{\MatterSpace}{\Psi}

\newcommand{\Car}{\mathrm{car}}
\newcommand{\Cur}{\mathrm{cur}}
\newcommand{\Hom}{\mathrm{hom}}

\newcommand{\ForSpace}[1]{\mathcal I_{#1}^{\mathrm{for}}}
\newcommand{\PrimShell}{\mathcal S_{\mathrm{prim}}}

\newcommand{\Reduction}[1]{\mathcal R_{#1}}
\newcommand{\CurrentCarrier}[1]{\mathfrak C_{#1}^{\mathrm{cur}}}
\newcommand{\EqCarrier}[1]{\mathfrak C_{#1}^{\mathrm{eq}}}

\newcommand{\MetricOut}{\mathscr G_{\mathrm{out}}}
\newcommand{\MatterOut}{\mathscr T_{\mathrm{out}}}

\newcommand{\VisibleAffine}[1]{\widehat X_{#1}}
\newcommand{\DataSpace}[1]{\mathcal Y_{#1}^{\mathrm{obs}}}
\newcommand{\AmbientTarget}[1]{E_{#1}^{\mathrm{amb}}}

\newcommand{\SubSpace}[1]{\mathcal M_{#1}^{\mathrm{sub}}}
\newcommand{\SubBody}[1]{\mathcal B_{#1}^{\mathrm{sub}}}
\newcommand{\SubTarget}[1]{E_{#1}^{\mathrm{sub}}}
\newcommand{\SubPair}[1]{\mathfrak s_{#1}^{\mathrm{sub}}}

\newcommand{\LiftSpace}[1]{\mathcal L_{#1}^{\mathrm{lift}}}
\newcommand{\LiftBody}[1]{\mathcal B_{#1}^{\mathrm{lift}}}
\newcommand{\PrecSpace}[1]{\mathcal P_{#1}^{\mathrm{prec}}}
\newcommand{\PrecBody}[1]{\mathcal B_{#1}^{\mathrm{prec}}}
\newcommand{\OrigSpace}[1]{\mathcal O_{#1}^{\mathrm{orig}}}
\newcommand{\OrigBody}[1]{\mathcal B_{#1}^{\mathrm{orig}}}

\newcommand{\GroundSpace}[1]{\mathcal G_{#1}^{\mathrm{grd}}}
\newcommand{\GroundBody}[1]{\mathcal B_{#1}^{\mathrm{grd}}}
\newcommand{\GroundProj}[1]{\Pi_{#1}^{\mathrm{grd}}}
\newcommand{\GroundOrigMin}[1]{\mathfrak f_{#1}^{\mathrm{grd,orig}}}
\newcommand{\GroundPrecMin}[1]{\mathfrak f_{#1}^{\mathrm{grd,prec}}}
\newcommand{\GroundLiftMin}[1]{\mathfrak f_{#1}^{\mathrm{grd,lift}}}
\newcommand{\GroundSubMin}[1]{\mathfrak m_{#1}^{\mathrm{grd}}}
\newcommand{\GroundSection}[1]{\Gamma_{#1}^{\mathrm{grd}}}
\newcommand{\GroundFunctional}[1]{\mathscr W_{#1}^{\mathrm{grd}}}
\newcommand{\GroundData}[1]{\Lambda_{#1}^{\mathrm{grd,dat}}}
\newcommand{\GroundShellMap}[1]{\Lambda_{#1}^{\mathrm{grd,sh}}}
\newcommand{\GroundShellSet}[1]{\mathcal Z_{#1}^{\mathrm{grd}}}
\newcommand{\GroundSym}[1]{\mathbb G_{#1}^{\mathrm{grd}}}
\newcommand{\GroundTime}[1]{\vartheta_{#1}^{\mathrm{grd}}}
\newcommand{\GroundChart}[1]{\Xi_{#1}^{\mathrm{grd}}}
\newcommand{\GroundSupport}[1]{\Theta_{#1}^{\mathrm{grd,sup}}}
\newcommand{\GroundSupportShell}[1]{\mathcal Z_{#1}^{\mathrm{grd,sup}}}
\newcommand{\GroundZeroCore}[1]{\mathcal Z_{#1}^{0,\mathrm{grd}}}
\newcommand{\GroundEqChart}[1]{\Psi_{#1}^{\mathrm{grd,sup}}}
\newcommand{\GroundZeroChart}[1]{\Psi_{#1}^{0,\mathrm{grd}}}
\newcommand{\GroundSplit}[1]{\Phi_{#1}^{\mathrm{grd}}}
\newcommand{\GroundCharge}[1]{\mathcal Q_{#1}^{\mathrm{grd}}}
\newcommand{\GroundResidual}[1]{\mathcal R_{#1}^{\mathrm{grd}}}
\newcommand{\GroundEmbed}[1]{\zeta_{#1}^{\mathrm{grd}}}
\newcommand{\GroundKernelCh}[1]{\widetilde K_{#1}^{0,\mathrm{grd}}}
\newcommand{\GroundStateComp}[1]{P_{#1}^{\mathrm{grd,co}}}

\newcommand{\FoundSpace}[1]{\mathcal F_{#1}^{\mathrm{fdn}}}
\newcommand{\FoundBody}[1]{\mathcal B_{#1}^{\mathrm{fdn}}}
\newcommand{\FoundProj}[1]{\Pi_{#1}^{\mathrm{fdn}}}
\newcommand{\FoundGroundMin}[1]{\mathfrak f_{#1}^{\mathrm{fdn,grd}}}
\newcommand{\FoundOrigMin}[1]{\mathfrak f_{#1}^{\mathrm{fdn,orig}}}
\newcommand{\FoundPrecMin}[1]{\mathfrak f_{#1}^{\mathrm{fdn,prec}}}
\newcommand{\FoundLiftMin}[1]{\mathfrak f_{#1}^{\mathrm{fdn,lift}}}
\newcommand{\FoundSubMin}[1]{\mathfrak m_{#1}^{\mathrm{fdn}}}
\newcommand{\FoundSection}[1]{\Gamma_{#1}^{\mathrm{fdn}}}
\newcommand{\FoundFunctional}[1]{\mathscr W_{#1}^{\mathrm{fdn}}}
\newcommand{\FoundData}[1]{\Lambda_{#1}^{\mathrm{fdn,dat}}}
\newcommand{\FoundShellMap}[1]{\Lambda_{#1}^{\mathrm{fdn,sh}}}
\newcommand{\FoundShellSet}[1]{\mathcal Z_{#1}^{\mathrm{fdn}}}
\newcommand{\FoundSym}[1]{\mathbb G_{#1}^{\mathrm{fdn}}}
\newcommand{\FoundTime}[1]{\vartheta_{#1}^{\mathrm{fdn}}}
\newcommand{\FoundChart}[1]{\Xi_{#1}^{\mathrm{fdn}}}
\newcommand{\FoundSupport}[1]{\Theta_{#1}^{\mathrm{fdn,sup}}}
\newcommand{\FoundSupportShell}[1]{\mathcal Z_{#1}^{\mathrm{fdn,sup}}}
\newcommand{\FoundZeroCore}[1]{\mathcal Z_{#1}^{0,\mathrm{fdn}}}
\newcommand{\FoundEqChart}[1]{\Psi_{#1}^{\mathrm{fdn,sup}}}
\newcommand{\FoundZeroChart}[1]{\Psi_{#1}^{0,\mathrm{fdn}}}
\newcommand{\FoundSplit}[1]{\Phi_{#1}^{\mathrm{fdn}}}
\newcommand{\FoundCharge}[1]{\mathcal Q_{#1}^{\mathrm{fdn}}}
\newcommand{\FoundResidual}[1]{\mathcal R_{#1}^{\mathrm{fdn}}}
\newcommand{\FoundEmbed}[1]{\zeta_{#1}^{\mathrm{fdn}}}
\newcommand{\FoundKernelCh}[1]{\widetilde K_{#1}^{0,\mathrm{fdn}}}
\newcommand{\FoundStateComp}[1]{P_{#1}^{\mathrm{fdn,co}}}

\newcommand{\ResSpace}[1]{\mathcal U_{#1}^{\mathrm{sup,res}}}
\newcommand{\ResBody}[1]{\mathcal B_{#1}^{\mathrm{sup,res}}}
\newcommand{\ResTarget}[1]{S_{#1}^{\mathrm{sup,res}}}
\newcommand{\ResPair}[1]{\mathfrak m_{#1}^{\mathrm{sup,res}}}
\newcommand{\ResSection}[1]{\Theta_{#1}^{\mathrm{sup,res}}}
\newcommand{\ResData}[1]{\Lambda_{#1}^{\mathrm{sup,dat}}}
\newcommand{\ResShellMap}[1]{\Lambda_{#1}^{\mathrm{sup,sh}}}
\newcommand{\ResShellSet}[1]{\mathcal Z_{#1}^{\mathrm{sup,res}}}
\newcommand{\SeedHidden}[1]{\mathcal H_{#1}^{\mathrm{sup}}}
\newcommand{\SeedHiddenBody}[1]{D_{#1}^{\mathrm{sup}}}
\newcommand{\HiddenChart}[1]{\Psi_{#1}^{\mathrm{hid}}}
\newcommand{\ZeroBody}[1]{\mathcal B_{#1}^{0,\mathrm{sup}}}

\newcommand{\EqnSpace}[1]{\mathcal U_{#1}^{\mathrm{eqn}}}
\newcommand{\EqnTarget}[1]{Q_{#1}^{\mathrm{eqn}}}
\newcommand{\EqnPair}[1]{\mathfrak b_{#1}^{\mathrm{eqn}}}
\newcommand{\EqnData}[1]{\Lambda_{#1}^{\mathrm{eqn,dat}}}
\newcommand{\EqnShell}[1]{\Lambda_{#1}^{\mathrm{eqn,sh}}}
\newcommand{\EqnHidden}[1]{H_{#1}^{\mathrm{eqn}}}
\newcommand{\EqnZero}[1]{V_{#1}^{\mathrm{eqn},0}}
\newcommand{\EqnBody}[1]{\mathcal B_{#1}^{\mathrm{eqn}}}

\newcommand{\FoundCore}[1]{\mathfrak F_{#1}^{\mathrm{fdn,core}}}
\newcommand{\CoreMapGround}[1]{\mathcal E_{#1}^{\mathrm{grd}}}
\newcommand{\CoreMapOrig}[1]{\mathcal E_{#1}^{\mathrm{orig}}}
\newcommand{\CoreMapPrec}[1]{\mathcal E_{#1}^{\mathrm{prec}}}
\newcommand{\CoreMapLift}[1]{\mathcal E_{#1}^{\mathrm{lift}}}
\newcommand{\CoreMapSub}[1]{\mathcal E_{#1}^{\mathrm{sub}}}
\newcommand{\CoreMapShell}[1]{\mathcal E_{#1}^{\mathrm{sh}}}
\newcommand{\CoreMapSupport}[1]{\mathcal E_{#1}^{\mathrm{sup}}}
\newcommand{\CoreMapZero}[1]{\mathcal E_{#1}^{0}}

\newtheorem{definition}{Definition}
\newtheorem{lemma}{Lemma}
\newtheorem{proposition}{Proposition}
\newtheorem{remark}{Remark}
\newtheorem{corollary}{Corollary}
\newtheorem{theorem}{Theorem}

\title{\textbf{The \TheoryName{}}\\[0.4em]
\large Primitive-Reservoir Observer-Reduction\\
\large Ground-Generating Substrate Foundation Dynamics\\
\large Necessity / Uniqueness Theorem}
\author{}
\date{}

\begin{document}

\maketitle

\begin{abstract}
The newly recorded theorem deriving origin-generating substrate ground dynamics from ground-generating substrate foundation dynamics moved the honest burden to the foundation level, but under the current routing safeguard it did not by itself become the benchmark-facing frontier. That theorem is still an origin-side relay: it introduces a deeper foundation layer while preserving exactly the same one-metric observer package. The next benchmark-facing move therefore cannot honestly be another shell-, lift-, support-, reservoir-, or ground-level restatement, another reopening of stationary reduction, stationary Hessians, spectral generators, or selector mechanics, or any stronger promotion language. The correct sharper task is to remove the remaining freedom inside the foundation layer itself.

The present note takes exactly that route. It does not descend to a still deeper substrate class. Instead, for each sector it isolates the benchmark-relevant \emph{fiber-minimizing exact-shell core} of a ground-generating substrate foundation presentation: the foundation-to-ground projection and the five recorded fiber-minimizer images, the exact foundation shell and its observer-data and shell-response readouts, the shell chart to the same reservoir body, the support shell and zero core, the zero/hidden split in the same equation variables, the hidden charge and exact zero-response realization, and the fixed-data solvability and coercive control carried on that core. The theorem then proves that every foundation-faithful presentation inducing the fixed ground package determines exactly the same foundation core up to canonical core-faithful equivalence. In that precise sense, no additional benchmark-relevant foundation-level freedom remains on the active line once the foundation layer is required to recover the recorded ground package.

The benchmark boundary remains unchanged throughout. The one operative metric is still exactly \(\CompletedMetric\), the ordinary matter route is unchanged, there is no second operative metric, and the low-energy matter law is not enlarged. The scope is therefore disciplined. This is a sharper foundation-level theorem inside the active derivational program of \TheoryName{}, not a completed first-principles derivation of gravity from final substrate microphysics. Its honest next burden is deeper again: derive or justify the foundation dynamics from still more primitive substrate structure, or prove a genuinely smaller theorem-level collapse strictly inside that foundation core.
\end{abstract}

\section{Introduction}

The active derivational program of \TheoryName{} now has a sharper next task than another deeper relay. The immediately preceding ground-from-foundation theorem already showed that the origin-generating substrate ground package is no longer the right disciplined stopping point on the active primitive-reservoir observer-reduction line. Once ground-generating substrate foundation dynamics are granted, the ground package is already their canonical projected exact-shell output. That result matters because it moves the burden below the ground layer. But it also sharpens what can count as the next benchmark-facing gain.

By the current routing rule, another same-output deeper-origin continuation does not become front-line progress merely by adding depth. It becomes front-line only if it changes benchmark-facing status by compression, necessity, uniqueness, or some comparably sharp theorem. The honest next task is therefore not to restate ground, origin, precursor, lift, shell-restriction, or support data once more. It is to ask whether the remaining freedom inside the foundation layer itself can already be removed while preserving the same benchmark one-metric package and the same claim boundary.

That is the point of the present note. The theorem proved here does not claim that final substrate microphysics has been found. Its gain is narrower and sharper. It shows that, at the current foundation level, the part of the foundation package that actually matters for the active one-metric derivational line is no longer a free choice. Once one requires a foundation presentation to induce exactly the fixed ground package, the benchmark-relevant foundation core is forced up to canonical equivalence.

\paragraph{Framework and claim boundary.}
\TheoryProgram{} denotes the broader \Aether{} / \Aflow{} framework built on the claims that the \Aether{} is the underlying four-dimensional substrate of reality, the \Aflow{} is its intrinsic ordered motion, observed three-dimensional space is the local experiential slice of that deeper substrate, S-time is the experienced order of change arising from matter, light, and the \Aflow{}, and observed expansion is the three-dimensional appearance of deeper four-dimensional ordered motion. Within that framework, \TheoryName{} denotes the exact-closure theory in which the effective gravitational dynamics are adopted to be exactly Einsteinian with universal matter coupling. In the active sequence, \emph{adoption} means use of that established relativistic dynamics without claiming substrate derivation, while \emph{derivation} is reserved for a first-principles recovery from explicit substrate variables.

The present note stays strictly inside that boundary. It does not alter the benchmark exact-closure package. It does not introduce a second operative metric or a new low-energy matter law. It only proves that the current foundation layer is already rigid at benchmark scope once it is required to reproduce the fixed ground package.

\section{Fixed Primitive Continuation Data}

\begin{definition}[Recorded primitive continuation data]
\label{def:fixed}
Fix the following continuation data from the active primitive-reservoir observer-reduction line.
\begin{enumerate}
    \item For each sector label \(\sigma\in\{\Car,\Cur,\Hom\}\), a primitive forcing shell
    \[
    \ForSpace{\sigma},
    \]
    a visible reduction map
    \[
    \Reduction{\sigma}:\ForSpace{\sigma}\to X_{\sigma},
    \]
    and shell-level current and equilibrium carriers
    \[
    \CurrentCarrier{\sigma}:\ForSpace{\sigma}\to Z_{\sigma}^{\mathrm{cur}},
    \qquad
    \EqCarrier{\sigma}:\ForSpace{\sigma}\to Z_{\sigma}^{\mathrm{eq}}.
    \]
    \item The product shell
    \[
    \PrimShell
    \equiv
    \ForSpace{\Car}\times \ForSpace{\Cur}\times \ForSpace{\Hom}.
    \]
    \item The recorded metric and ordinary-matter outputs
    \[
    \MetricOut:\PrimShell\to\{\text{observer metrics}\},
    \qquad
    \MatterOut:\PrimShell\times\MatterSpace\to\{\text{observer stress tensors}\},
    \]
    satisfying
    \begin{equation}
    \MetricOut(\mathbf w)=\CompletedMetric,
    \qquad
    \MatterOut(\mathbf w,\psi)
    =
    T_{\mu\nu}^{\mathrm{matter}}[\CompletedMetric,\psi]
    \label{eq:fixedoutputs}
    \end{equation}
    for every \(\mathbf w\in\PrimShell\) and every matter configuration \(\psi\in\MatterSpace\).
\end{enumerate}
\end{definition}

\begin{remark}
Definition \ref{def:fixed} is not reopened here. The primitive continuation data, the one operative metric, and the ordinary matter route remain fixed. The present note only addresses the remaining freedom inside the already recorded foundation layer that now sits below the ground theorem on the active line.
\end{remark}

\section{Recorded Ground Target}

\begin{definition}[Recorded origin-generating substrate ground package]
\label{def:ground}
Fix \(\sigma\in\{\Car,\Cur,\Hom\}\). The fixed ground package on the active line consists of the following data.
\begin{enumerate}
    \item A finite-dimensional Euclidean ground state space \(\GroundSpace{\sigma}\), a compact convex ground body \(\GroundBody{\sigma}\subset\GroundSpace{\sigma}\) with nonempty interior, and fixed lower bodies
    \[
    \OrigBody{\sigma},\qquad
    \PrecBody{\sigma},\qquad
    \LiftBody{\sigma},\qquad
    \SubBody{\sigma},
    \]
    together with an affine surjection
    \[
    \GroundProj{\sigma}:\GroundSpace{\sigma}\to\OrigSpace{\sigma}
    \]
    and unique ground fiber minimizers
    \[
    \GroundOrigMin{\sigma},\quad
    \GroundPrecMin{\sigma},\quad
    \GroundLiftMin{\sigma},\quad
    \GroundSubMin{\sigma}
    \]
    on the already recorded origin, precursor, lift, and substrate fibers of the active lower line.
    \item The same finite-dimensional Euclidean substrate restriction target \(\SubTarget{\sigma}\), the same positive-definite symmetric substrate pairing
    \[
    \SubPair{\sigma}:\SubTarget{\sigma}\times\SubTarget{\sigma}\to\mathbb R,
    \]
    a smooth ground restriction section
    \[
    \GroundSection{\sigma}:\GroundSpace{\sigma}\to\SubTarget{\sigma},
    \]
    the ground functional
    \[
    \GroundFunctional{\sigma}(g)
    \equiv
    \SubPair{\sigma}\bigl(\GroundSection{\sigma}(g),\GroundSection{\sigma}(g)\bigr),
    \]
    and the exact ground shell
    \[
    \GroundShellSet{\sigma}
    \equiv
    \bigl(\GroundSection{\sigma}\bigr)^{-1}(0)\cap\GroundBody{\sigma}.
    \]
    \item Affine observer-data and shell-response readouts
    \[
    \GroundData{\sigma}:\GroundSpace{\sigma}\to\DataSpace{\sigma},
    \qquad
    \GroundShellMap{\sigma}:\GroundSpace{\sigma}\to\AmbientTarget{\sigma},
    \]
    together with a compact symmetry group \(\GroundSym{\sigma}\) and an involution \(\GroundTime{\sigma}\) preserving the displayed data and bodies.
    \item A Euclidean shell chart
    \[
    \GroundChart{\sigma}:\GroundShellSet{\sigma}\to\ResSpace{\sigma}
    \]
    whose image
    \[
    \ResBody{\sigma}
    \equiv
    \GroundChart{\sigma}\bigl(\GroundShellSet{\sigma}\bigr)
    \]
    is compact convex with nonempty interior, affine readouts
    \[
    \ResData{\sigma}:\ResSpace{\sigma}\to\DataSpace{\sigma},
    \qquad
    \ResShellMap{\sigma}:\ResSpace{\sigma}\to\AmbientTarget{\sigma},
    \]
    a smooth support section
    \[
    \ResSection{\sigma}:\ResSpace{\sigma}\to\ResTarget{\sigma},
    \]
    a shell-internal support recorder
    \[
    \GroundSupport{\sigma}:\GroundShellSet{\sigma}\to\ResTarget{\sigma},
    \]
    and the exact ground support shell
    \[
    \GroundSupportShell{\sigma}
    \equiv
    \bigl(\GroundSupport{\sigma}\bigr)^{-1}(0)\cap\GroundShellSet{\sigma}.
    \]
    The support section satisfies
    \[
    \ResShellSet{\sigma}
    \equiv
    \bigl(\ResSection{\sigma}\bigr)^{-1}(0)\cap\ResBody{\sigma}
    =
    \GroundChart{\sigma}\bigl(\GroundSupportShell{\sigma}\bigr),
    \]
    together with the normal estimate already recorded on the active line.
    \item A closed embedded ground support zero core
    \[
    \GroundZeroCore{\sigma}\subset\GroundSupportShell{\sigma},
    \]
    a hidden-seed space \(\SeedHidden{\sigma}\) with compact convex body \(\SeedHiddenBody{\sigma}\) containing \(0\) in its interior, a support-shell chart
    \[
    \GroundEqChart{\sigma}:\GroundSupportShell{\sigma}\to\EqnSpace{\sigma},
    \]
    an affine zero chart
    \[
    \GroundZeroChart{\sigma}:\GroundZeroCore{\sigma}\to\EqnZero{\sigma},
    \]
    a linear isometric embedding
    \[
    \HiddenChart{\sigma}:\SeedHidden{\sigma}\to\EqnHidden{\sigma},
    \]
    the orthogonal splitting
    \[
    \EqnSpace{\sigma}
    =
    \EqnZero{\sigma}\oplus\EqnHidden{\sigma},
    \]
    and a diffeomorphism
    \[
    \GroundSplit{\sigma}:
    \GroundZeroCore{\sigma}\times\SeedHiddenBody{\sigma}
    \xrightarrow{\cong}
    \GroundSupportShell{\sigma}
    \]
    satisfying the chart split
    \[
    \GroundEqChart{\sigma}\bigl(\GroundSplit{\sigma}(z,\eta)\bigr)
    =
    \GroundZeroChart{\sigma}(z)+\HiddenChart{\sigma}(\eta).
    \]
    There are affine equation readouts
    \[
    \EqnData{\sigma}:\EqnSpace{\sigma}\to\DataSpace{\sigma},
    \qquad
    \EqnShell{\sigma}:\EqnSpace{\sigma}\to\AmbientTarget{\sigma}
    \]
    compatible with the ground support-shell chart.
    \item A linear hidden charge
    \[
    \GroundCharge{\sigma}:\SeedHidden{\sigma}\to\EqnTarget{\sigma}
    \]
    with positive hidden gap, the hidden recorder
    \[
    \GroundResidual{\sigma}\bigl(\GroundSplit{\sigma}(z,\eta)\bigr)
    \equiv
    \GroundCharge{\sigma}(\eta),
    \]
    and a smooth observer-compatible exact zero-response realization
    \[
    \GroundEmbed{\sigma}:\ForSpace{\sigma}\hookrightarrow\GroundZeroCore{\sigma}
    \]
    recovering the fixed visible, current, and equilibrium shell data.
    \item Fixed-data solvability on \(\GroundZeroCore{\sigma}\) and coercive control on data-invisible directions, recorded through the charted kernel \(\GroundKernelCh{\sigma}\) and orthogonal projector \(\GroundStateComp{\sigma}\).
\end{enumerate}
\end{definition}

\begin{remark}
Definition \ref{def:ground} is the fixed target now carried by the recorded ground-from-foundation theorem. The present note does not introduce a different ground package. It asks whether the remaining foundation-level freedom above that package can already be removed at theorem level.
\end{remark}

\section{Foundation Presentations and Their Core}

\begin{definition}[Ground-generating substrate foundation dynamics]
\label{def:foundation}
Fix \(\sigma\in\{\Car,\Cur,\Hom\}\). A ground-generating substrate foundation presentation for sector \(\sigma\) consists of the following data.
\begin{enumerate}
    \item A finite-dimensional Euclidean foundation state space \(\FoundSpace{\sigma}\), a compact convex foundation body \(\FoundBody{\sigma}\subset\FoundSpace{\sigma}\) with nonempty interior, an affine surjection
    \[
    \FoundProj{\sigma}:\FoundSpace{\sigma}\to\GroundSpace{\sigma},
    \]
    and unique foundation fiber minimizers
    \[
    \FoundGroundMin{\sigma},\quad
    \FoundOrigMin{\sigma},\quad
    \FoundPrecMin{\sigma},\quad
    \FoundLiftMin{\sigma},\quad
    \FoundSubMin{\sigma}
    \]
    on the fixed ground, origin, precursor, lift, and substrate fibers of the active line.
    \item The same substrate target \(\SubTarget{\sigma}\), the same pairing \(\SubPair{\sigma}\), a smooth foundation section
    \[
    \FoundSection{\sigma}:\FoundSpace{\sigma}\to\SubTarget{\sigma},
    \]
    the foundation functional
    \[
    \FoundFunctional{\sigma}(f)
    \equiv
    \SubPair{\sigma}\bigl(\FoundSection{\sigma}(f),\FoundSection{\sigma}(f)\bigr),
    \]
    and the exact foundation shell
    \[
    \FoundShellSet{\sigma}
    \equiv
    \bigl(\FoundSection{\sigma}\bigr)^{-1}(0)\cap\FoundBody{\sigma}.
    \]
    \item Affine foundation observer-data and shell-response readouts
    \[
    \FoundData{\sigma}:\FoundSpace{\sigma}\to\DataSpace{\sigma},
    \qquad
    \FoundShellMap{\sigma}:\FoundSpace{\sigma}\to\AmbientTarget{\sigma},
    \]
    together with a compact symmetry group \(\FoundSym{\sigma}\) and an involution \(\FoundTime{\sigma}\) preserving all displayed foundation data.
    \item A Euclidean foundation shell chart
    \[
    \FoundChart{\sigma}:\FoundShellSet{\sigma}\to\ResSpace{\sigma},
    \]
    a shell-internal support recorder
    \[
    \FoundSupport{\sigma}:\FoundShellSet{\sigma}\to\ResTarget{\sigma},
    \]
    the exact foundation support shell
    \[
    \FoundSupportShell{\sigma}
    \equiv
    \bigl(\FoundSupport{\sigma}\bigr)^{-1}(0)\cap\FoundShellSet{\sigma},
    \]
    a closed embedded foundation zero core
    \[
    \FoundZeroCore{\sigma}\subset\FoundSupportShell{\sigma},
    \]
    a support-shell chart
    \[
    \FoundEqChart{\sigma}:\FoundSupportShell{\sigma}\to\EqnSpace{\sigma},
    \]
    a zero chart
    \[
    \FoundZeroChart{\sigma}:\FoundZeroCore{\sigma}\to\EqnZero{\sigma},
    \]
    and a diffeomorphism
    \[
    \FoundSplit{\sigma}:
    \FoundZeroCore{\sigma}\times\SeedHiddenBody{\sigma}
    \xrightarrow{\cong}
    \FoundSupportShell{\sigma}
    \]
    obeying the same zero/hidden split in the same equation variables.
    \item A hidden charge
    \[
    \FoundCharge{\sigma}:\SeedHidden{\sigma}\to\EqnTarget{\sigma}
    \]
    with positive hidden gap, the hidden recorder
    \[
    \FoundResidual{\sigma}\bigl(\FoundSplit{\sigma}(z,\eta)\bigr)
    \equiv
    \FoundCharge{\sigma}(\eta),
    \]
    and a smooth observer-compatible exact foundation zero-response realization
    \[
    \FoundEmbed{\sigma}:\ForSpace{\sigma}\hookrightarrow\FoundZeroCore{\sigma}
    \]
    recovering the fixed visible, current, and equilibrium shell data.
    \item Fixed-data solvability on \(\FoundZeroCore{\sigma}\) and coercive control on foundation data-invisible directions, recorded through the charted kernel \(\FoundKernelCh{\sigma}\) and orthogonal projector \(\FoundStateComp{\sigma}\).
\end{enumerate}
\end{definition}

\begin{definition}[Foundation-faithful presentation]
\label{def:faithful}
A foundation presentation in the sense of Definition \ref{def:foundation} is \emph{foundation-faithful} if its projected exact-shell transport recovers the fixed ground package of Definition \ref{def:ground} exactly, sector by sector, while preserving the benchmark outputs \eqref{eq:fixedoutputs}. Equivalently, the foundation presentation is required to induce precisely the recorded ground body, the same ground fiber-minimizer hierarchy, the same ground shell and support transport, the same zero/hidden split in the same equation variables, the same hidden charge data, the same exact zero-response realization, and the same solvability/coercive statements, with no second operative metric and no enlarged low-energy matter law.
\end{definition}

\begin{definition}[Fiber-minimizing exact-shell foundation core]
\label{def:core}
For a foundation-faithful presentation, the \emph{fiber-minimizing exact-shell foundation core} \(\FoundCore{\sigma}\) is the tuple consisting of:
\begin{enumerate}
    \item the projection \(\FoundProj{\sigma}\) together with the images of the five fiber minimizers \(\FoundGroundMin{\sigma}\), \(\FoundOrigMin{\sigma}\), \(\FoundPrecMin{\sigma}\), \(\FoundLiftMin{\sigma}\), and \(\FoundSubMin{\sigma}\);
    \item the exact foundation shell \(\FoundShellSet{\sigma}\), the foundation readouts \(\FoundData{\sigma}\) and \(\FoundShellMap{\sigma}\), the shell chart \(\FoundChart{\sigma}\), and the support recorder \(\FoundSupport{\sigma}\);
    \item the support shell \(\FoundSupportShell{\sigma}\), the zero core \(\FoundZeroCore{\sigma}\), the support-shell chart \(\FoundEqChart{\sigma}\), the zero chart \(\FoundZeroChart{\sigma}\), and the split diffeomorphism \(\FoundSplit{\sigma}\);
    \item the hidden charge \(\FoundCharge{\sigma}\), the hidden recorder \(\FoundResidual{\sigma}\), the exact zero-response realization \(\FoundEmbed{\sigma}\), and the solvability/coercive statements carried by \(\FoundKernelCh{\sigma}\) and \(\FoundStateComp{\sigma}\).
\end{enumerate}
\end{definition}

\begin{definition}[Core-faithful equivalence]
\label{def:equiv}
Let \({}^{(1)}\!\FoundCore{\sigma}\) and \({}^{(2)}\!\FoundCore{\sigma}\) be the foundation cores of two foundation-faithful presentations of the same fixed ground package. Their canonical comparison maps are defined by the lower data that both cores induce:
\begin{align}
\CoreMapGround{\sigma}\bigl({}^{(1)}\!\FoundGroundMin{\sigma}(g)\bigr)
&\equiv
{}^{(2)}\!\FoundGroundMin{\sigma}(g),
\label{eq:eqgrd}
\\
\CoreMapOrig{\sigma}\bigl({}^{(1)}\!\FoundOrigMin{\sigma}(o)\bigr)
&\equiv
{}^{(2)}\!\FoundOrigMin{\sigma}(o),
\label{eq:eqorig}
\\
\CoreMapPrec{\sigma}\bigl({}^{(1)}\!\FoundPrecMin{\sigma}(p)\bigr)
&\equiv
{}^{(2)}\!\FoundPrecMin{\sigma}(p),
\label{eq:eqprec}
\\
\CoreMapLift{\sigma}\bigl({}^{(1)}\!\FoundLiftMin{\sigma}(\ell)\bigr)
&\equiv
{}^{(2)}\!\FoundLiftMin{\sigma}(\ell),
\label{eq:eqlift}
\\
\CoreMapSub{\sigma}\bigl({}^{(1)}\!\FoundSubMin{\sigma}(m)\bigr)
&\equiv
{}^{(2)}\!\FoundSubMin{\sigma}(m),
\label{eq:eqsub}
\\
\CoreMapShell{\sigma}
&\equiv
\bigl({}^{(2)}\!\FoundProj{\sigma}\big|_{{}^{(2)}\!\FoundShellSet{\sigma}}\bigr)^{-1}
\circ
{}^{(1)}\!\FoundProj{\sigma}\big|_{{}^{(1)}\!\FoundShellSet{\sigma}},
\label{eq:eqshell}
\\
\CoreMapSupport{\sigma}
&\equiv
\bigl({}^{(2)}\!\FoundProj{\sigma}\big|_{{}^{(2)}\!\FoundSupportShell{\sigma}}\bigr)^{-1}
\circ
{}^{(1)}\!\FoundProj{\sigma}\big|_{{}^{(1)}\!\FoundSupportShell{\sigma}},
\label{eq:eqsupport}
\\
\CoreMapZero{\sigma}
&\equiv
\bigl({}^{(2)}\!\FoundProj{\sigma}\big|_{{}^{(2)}\!\FoundZeroCore{\sigma}}\bigr)^{-1}
\circ
{}^{(1)}\!\FoundProj{\sigma}\big|_{{}^{(1)}\!\FoundZeroCore{\sigma}}.
\label{eq:eqzero}
\end{align}
The two cores are \emph{core-faithfully equivalent} if, under the maps \eqref{eq:eqgrd}--\eqref{eq:eqzero}, every displayed structure in Definition \ref{def:core} is transported identically.
\end{definition}

\begin{remark}
Definition \ref{def:equiv} is intentionally narrower than full off-shell identity of the ambient foundation body. The active burden is not to classify every possible foundation-thickening outside the fiber-minimizing exact-shell core. It is to remove the benchmark-relevant foundation freedom still visible on the active one-metric line.
\end{remark}

\section{Necessity of the Foundation Core}

\begin{proposition}[The fiber-minimizer hierarchy is necessary]
\label{prop:minimizers}
Every foundation-faithful presentation determines the same ground-, origin-, precursor-, lift-, and substrate-fiber minimizer images inside the foundation layer, hence the maps \eqref{eq:eqgrd}--\eqref{eq:eqsub} are well defined.
\end{proposition}

\begin{proof}
By Definition \ref{def:faithful}, the foundation presentation must recover exactly the fixed ground package of Definition \ref{def:ground}. In particular, it must recover the same ground body and the same lower bodies together with the same unique minimizers on the corresponding recorded fibers. Since those lower fibers are fixed continuation data, the image of each foundation minimizer is keyed uniquely by the lower point it lifts. The comparison maps \eqref{eq:eqgrd}--\eqref{eq:eqsub} therefore depend only on the shared lower datum and are canonical.
\end{proof}

\begin{proposition}[Exact-shell transport data are necessary]
\label{prop:shell}
Every foundation-faithful presentation determines the same exact-shell transport data at benchmark scope. More precisely, the projection \(\FoundProj{\sigma}\) restricts to diffeomorphisms from the foundation shell, support shell, and zero core onto the corresponding fixed ground loci, and the shell chart, support recorder, support-shell chart, zero chart, zero/hidden split, and equation readouts are thereby forced up to the canonical comparison maps \eqref{eq:eqshell}--\eqref{eq:eqzero}.
\end{proposition}

\begin{proof}
If a foundation presentation is faithful, then by definition its projected exact-shell transport recovers precisely the shell-side data already fixed in Definition \ref{def:ground}. The foundation shell must therefore project diffeomorphically to the recorded ground shell, otherwise the fixed ground shell would not be recovered as exact-shell image. The same argument applies one step further to the support shell and then to the zero core. Once those three projected loci are fixed, the shell chart to \(\ResBody{\sigma}\), the support recorder, the support-shell chart to \(\EqnSpace{\sigma}\), the zero chart to \(\EqnZero{\sigma}\), and the zero/hidden split in the same equation variables cannot vary independently: changing any of them would change the projected ground shell transport or the fixed equation-side coordinates already recorded on the active line. Hence all these constituents are forced up to transport by \eqref{eq:eqshell}--\eqref{eq:eqzero}.
\end{proof}

\begin{proposition}[Hidden charge, zero-response realization, solvability, and coercivity are necessary]
\label{prop:hidden}
Every foundation-faithful presentation determines, up to the canonical comparison maps, the same hidden charge data, the same exact zero-response realization, the same fixed-data solvability statement, and the same coercive control on data-invisible directions.
\end{proposition}

\begin{proof}
The fixed ground package already records one hidden-seed space \(\SeedHidden{\sigma}\), one hidden embedding \(\HiddenChart{\sigma}\), one hidden charge \(\GroundCharge{\sigma}\) with positive gap, one exact zero-response realization \(\GroundEmbed{\sigma}\), and one solvability/coercive package on the zero core. A foundation-faithful presentation must project to that same ground package. Because the support-shell split and the zero/hidden coordinates are already fixed by Proposition \ref{prop:shell}, the hidden sector cannot be changed without changing the induced ground recorder. Likewise, the exact zero-response realization must project to the same fixed visible/current/equilibrium shell data, so its image on the foundation zero core is forced by the same ground compatibility requirement. The solvability and coercive statements are properties of this same projected zero-response core; if they differed at foundation scope, the induced ground statements would differ as well. Therefore those constituents are necessary up to the canonical comparison maps.
\end{proof}

\begin{proposition}[The benchmark package remains fixed]
\label{prop:benchmark}
Every foundation-faithful presentation preserves the same benchmark one-metric package \eqref{eq:fixedoutputs}. In particular, the operative metric remains exactly \(\CompletedMetric\), the ordinary matter route is unchanged, there is no second operative metric, and the low-energy matter law is not enlarged.
\end{proposition}

\begin{proof}
This is part of Definition \ref{def:faithful}. The present note removes only foundation-level freedom internal to the derivational line. It does not alter the recorded primitive outputs.
\end{proof}

\section{Main Theorem}

\begin{theorem}[Ground-generating substrate foundation dynamics necessity / uniqueness]
\label{thm:main}
Fix the primitive continuation data of Definition \ref{def:fixed} and the recorded ground package of Definition \ref{def:ground}. Then, for each sector \(\sigma\in\{\Car,\Cur,\Hom\}\), the fiber-minimizing exact-shell foundation core \(\FoundCore{\sigma}\) is necessary and unique at benchmark scope up to core-faithful equivalence. More precisely:
\begin{enumerate}
    \item every foundation-faithful presentation must determine the same fiber-minimizer hierarchy and the same exact-shell transport data described in Definition \ref{def:core};
    \item for any two foundation-faithful presentations, the canonical comparison maps \eqref{eq:eqgrd}--\eqref{eq:eqzero} are well defined and yield a unique core-faithful equivalence;
    \item under that equivalence, the foundation shell, support shell, zero core, shell/support/equation charts, hidden charge, exact zero-response realization, fixed-data solvability statement, and coercive control all coincide;
    \item consequently no additional benchmark-relevant foundation-level freedom remains on the active one-metric line beyond core-faithful relabeling.
\end{enumerate}
\end{theorem}

\begin{proof}
Item (1) is exactly the combined content of Propositions \ref{prop:minimizers}, \ref{prop:shell}, and \ref{prop:hidden}. Those propositions show that no foundation-faithful presentation may vary the fiber-minimizing exact-shell core while still inducing the same fixed ground package.

For item (2), Proposition \ref{prop:minimizers} gives the canonical comparison maps on the five minimizer images, while Proposition \ref{prop:shell} gives the canonical comparison maps on the shell, support shell, and zero core. Because each of those loci projects to the same fixed lower datum, the maps \eqref{eq:eqgrd}--\eqref{eq:eqzero} are unique.

For item (3), Proposition \ref{prop:shell} shows that the shell chart, support recorder, support-shell chart, zero chart, split, and equation readouts are transported identically under the canonical comparison maps. Proposition \ref{prop:hidden} then shows that the hidden charge, exact zero-response realization, solvability statement, and coercive control are likewise transported identically. Hence the entire foundation core agrees under the unique core-faithful equivalence.

Item (4) is the project-level consequence of the previous items. Any remaining off-shell freedom outside the fiber-minimizing exact-shell core is not benchmark-relevant on the current line, because it does not change the fixed ground package or the benchmark outputs. Therefore the current honest foundation burden has been removed in the only sense that matters for the active one-metric derivational program.
\end{proof}

\begin{corollary}[No remaining benchmark-relevant foundation choice]
\label{cor:nofreedom}
At current scope there is no second benchmark-relevant foundation presentation of the recorded ground package other than a core-faithful relabeling of \(\FoundCore{\sigma}\) in each sector.
\end{corollary}

\begin{proof}
If a second benchmark-relevant foundation presentation existed, it would either change some constituent proved necessary in Theorem \ref{thm:main}, contradicting item (1), or else differ only by the canonical identifications of item (2), in which case it would be merely a core-faithful relabeling.
\end{proof}

\begin{remark}
Corollary \ref{cor:nofreedom} is the precise benchmark-facing sense in which the present manuscript is stronger than another deeper-origin relay. The current foundation layer is no longer merely available because a still deeper theorem reproduces it. The benchmark-relevant foundation core is now forced at current scope.
\end{remark}

\section{Routing Consequence}

\begin{corollary}[What must be done next]
\label{cor:next}
After Theorem \ref{thm:main}, the next honest primary manuscript below the current frontier must do at least one of the following:
\begin{enumerate}
    \item derive or justify the ground-generating substrate foundation dynamics from still more primitive substrate structure;
    \item identify a genuinely smaller theorem-level core strictly inside \(\FoundCore{\sigma}\) and prove a sharper collapse to it;
    \item do either of the previous two while preserving the same benchmark one-metric package, the same ordinary matter route, and the same claim boundary.
\end{enumerate}
It is no longer enough merely to restate the current foundation layer or to insert another same-output relay above it.
\end{corollary}

\begin{proof}
Theorem \ref{thm:main} removes the benchmark-relevant freedom presently visible at the foundation level. The next honest burden therefore lies either below that level or in a strictly smaller internal compression of it. Another package-preserving restatement of the same foundation core would not change project status.
\end{proof}

\section{Conclusion}

The positive result of this note is not that a final substrate microphysics has already been found. Its gain is narrower and more honest. The recorded ground-from-foundation theorem had already shown that the current ground package is the canonical projected exact-shell output of a deeper foundation class. The present theorem then removes the remaining benchmark-relevant freedom inside that foundation layer itself. At current scope, the fiber-minimizing exact-shell foundation core is necessary and unique up to canonical core-faithful equivalence.

That conclusion is exactly the sharper theorem-level gain the project needed. The foundation-to-ground projection, the five fiber-minimizer images, the foundation shell and support transport, the zero/hidden split in the same equation variables, the hidden charge, the exact zero-response realization, and the solvability/coercive package are no longer merely one possible deeper presentation of the same downstream line. Once the presentation is required to induce the fixed ground package, those constituents are forced.

The scope remains disciplined. The benchmark package is unchanged: the one operative metric is still exactly \(\CompletedMetric\), the ordinary matter route is unchanged, there is no second operative metric, and the low-energy matter law is not enlarged. This remains a theorem inside the active derivational program of \TheoryName{}, not a basis for stronger promotion language. The next honest burden is deeper again: derive or justify the foundation dynamics from still more primitive substrate structure, or else prove a genuinely smaller collapse inside the current foundation core. Until then, adoption and derivation should continue to be kept sharply separated.

\end{document}
